\documentclass{thesis_utf8}

\usepackage{relsize}
\newcommand*{\defeq}{\stackrel{\text{def}}{=}}

\usepackage{enumitem}

\begin{document}
\maketitlepage{Погода Михайло Володимирович}
    {КМ-31М}
    {завідувач кафедри, д-р техн. наук, доцент Чертов~О.\,Р.}
    {професор Тарасенко~В.\,П.}
    {Пошук шаблонів у цифровому сигналі за допомогою поліномів Кунченка}
\annotation{Анотація}
% TODO: Написати анотацію
Робота складається з вступу, \total{chapter}  розділів та висновків і налічує \total{page} сторінок.
Містить \total{figurecount} ілюстративних матеріалів, \total{tablecount} таблиць, \total{appendnum} додатки та
посилається на \total{bibitemcount} літературних джерел.

\clearpage

\annotation{Abstract}
% TODO: Write abstract
The work consists of an introduction, \total{chapter} sections, includes conclusions and \total{page} pages.
Contains \total{figurecount} illustrative materials, \total{tablecount} tables, \total{appendnum} appendices and has
\total{bibitemcount} references.

% vim: spelllang=uk,en spell

\tableofcontents
\shortings{}
CC --- (з англ. Cross"=correlation) взаємо"=кореляція

SSD --- (з англ. Sum of Squared Differencies) сума квадратів відстаней

NCC --- (з англ. Normalised Cross"=correlation) нормалізована взаємо"=кореляція

NSSD --- (з англ. Normalised Sum of Squared Differencies) нормалізована сума квадратів відстаней

Корелянта --- Скалярний добуток у функціональному просторі.

% vim: spelllang=uk,en spell

\intro{}
Інтелектуальна обробка інформації вже дуже давно зарекомендувала себе в сучасному світі.
Людство генерує незчисленні обсяги інформації щосекунди.
Багато згенерованої інформації (поки що) не зберігається, а з того, що зберігається, багато ніяк не аналізується.

Але з поширенням інформаційних технологій та з розповсюдженням пристроїв, що можуть ефективно обробляти великі обсяги
інформації, знаходиться все більше сфер, де можливо отримати користь з раніш не оброблюваного інформаційного потоку.

Наприклад, аналізуючи зміни погодних умов упродовж дня, метеорологи мають змогу прогнозувати погоду наперед;
аналізуючи різні показники життєдіяльності людини, лікарі знаходять прояви хвороб ще до перших її симптомів.

Задачі пошуку окремих наперед визначених властивостей чи образів у деякому масиві даних набули великого поширення в
галузях, пов’язаних з обробкою цифрової інформації.

Співставлення з еталоном (template matching) широко застосовується в аналізі як одновимірних, так і двовимірних
сигналів~\cite{book4}.

Зазвичай, постановка такої задачі включає в себе дані для аналізу, деякий шаблон для пошуку та, можливо, область, на
якій потрібно здійснювати такий пошук.  При цьому потрібно як правильно обрати представлення вхідних даних, так і
визначити критерії, за якими відбуватиметься зіставлення шаблону з сигналом~\cite{book10}.

Існує розмаїття підходів та метрик, які можна використовувати для вирішення цієї задачі.

У даній роботі для цього пропонується використовувати поліноми Кунченка.

% vim: spelllang=uk,en spell

\chapter{Постановка задачі}

\todo[inline]{Описати постановку задачі}
% vim: spelllang=uk,en spell

\chapter{Огляд існуючих методів}

Оскільки задача пошуку шаблонів у сигналах була поставлена дуже давно, існує багато методів пошуку шаблону в сигналі.
Ці методи різняться за областю застосування, швидкодією, чутливістю до різних перетворень шаблону.

В цьому розділі будуть розглянуті основні існуючі методи:
\begin{itemize}
    \item Взаємокореляція;
    \item Нормалізована взаємокореляція;
    \item Сума квадратів відстаней
\end{itemize}

Буде проведений порівняльний аналіз методів й обґрунтована доцільність дослідження обраного методу.

\section{Опис предметної області}
    В цьому підрозділі будуть описані загальні терміни, що будуть необхідні при описі методів.

    Ми будемо використовувати в якості сигналу довільну функцію
    \[f:\:\mathbb{X} \rightarrow \mathbb{Y},\] де множина вихідних значень $\mathbb{Y} \in \mathbb{R}$, а множина
    вхідних значень $\mathbb{X} \in \mathbb{R}$ для одновимірного сигналу, $\mathbb{X} \in \mathbb{R}^2$ для двовимірного сигналу, тощо.

    Під цифровим сигналом будемо розуміти функцію $f^*$, що задана співвідношеннями:
    \begin{eqnarray*}
        f^*(x_0) &=& y_0\\
        f^*(x_1) &=& y_1\\
        \dots\\
        f^*(x_n) &=& y_n
    \end{eqnarray*}
    Такі функції ще називають таблично"=заданими, оскільки вони можуть бути записаними у вигляді таблиці.

    Наприклад, для цифрового звукового сигналу, множина $\{x_i\}$ представляє собою моменти часу, в яких визначені
    амплітуда звукового сигналу $\{y_i\}$.
    % TODO s/"=/"=/g
    Для цифрового чорно"=білого зображення з розмірами $600 \times 400$ пікселів, область визначення цифрового сигналу
    буде множина пар
    \[\{ (a, b) \}, a \in \{1,2,\dots,599,600\}, b \in \{1,2,\dots,399,400\}\]
    Областю значень сигналу може бути, наприклад, $\{0,1,2,\dots,254,255\}$, де значення $0$ має чорний піксель, а
    $255$ --- білий.

    З кожної функції $f(x)$ можна утворити безліч таблично"=заданих функцій, вибравши з області визначення цієї
    функції скінченну підмножину $\{x_i\}$ й співставивши цієї підмножині відповідні значення функції $\{f(x_i)\}$.
    Такий процес ще називається дискретизацією сигналу.
    Якщо в якості вхідного сигналу виступає сигнал, що залежить від часу (наприклад, аудіо"=сигнал), за умови, що $|
    x_i - x_{i-1}| = \delta$, то частотою дискретизації називають співвідношення $\frac{1\,c}{\delta}$.

    \subsection{Метод ковзаючого вікна}
    \label{ss:sliding-window}
        Оскільки задача пошуку шаблону в цифровому сигналі ставиться як пошук наперед визначеного шаблону (або його
        трансформацій) в дискретному сигналі, що має (набагато) більшу довжину, то в методах пошуку шаблонів
        використовується метод ковзаючого вікна.

        Цей метод полягає в розбитті вхідного сигналу на \invcommas{вікна} --- неперервні ділянки, кожна з яких має
        одну й ту саму довжину.
        Як правило, довжина вікна дорівнює довжині шаблону, а самі вікна накладаються одне на одного.
        Для кожного вікна шукається міра подібності вмісту вікна до шаблону.

        Таким чином будується функція $e(x)$, яка характеризує подібність вікна, що задається параметром $x$ до
        шаблону.
        З цієї функції можна оцінювати найвірогідніші позиції шаблону.

        Така оцінка, як правило, робиться шукаючи локальні екстремуми, які вище/нище якогось порогу $e_{0}$.

\section{Опис існуючих методів}
    \subsection{Взаємокореляція}
        Взаємокореляція (cross-correlation) являє собою міру подібності двох функцій (сигналів), одна з яких
        зсувається відносно іншої.

        Для дискретного сигналу взаємокореляція має вигляд:
        \begin{equation}
            CC(f,g)[ n ] \equiv (f \star g)[ n ] \defeq \sum\limits_{m=0}^{k} f^*[m] g[ m + n ]
        \end{equation}

        Для використання цього методу необхідно взяти в якості сигналу $f$ вхідний сигнал, $g$ --- шаблон, що
        шукається; в якості $k$ необхідно взяти довжину (кількість значень) шаблону $g$.

        Ідея полягає в тому, що сума добутків значень шаблону на вхідний сигнал буде тим більше, чим більше вхідний
        сигнал схожий на шаблон на обраному проміжку.

        Значним недоліком цього методу є залежність від амплітуди сигналу: якщо сигнал має на деякому проміжку
        амплітуду, значно більшу за середню амплітуду сигналу та шаблону, то значення взаємокореляції між сигналом та
        шаблоном на цьому проміжку буде значно більше, ніж значення взаємокореляції між шаблоном та шаблоном
        (автокореляція).

        Саме тому більш практичним вважається використання нормалізованої взаємокореляції.
    \subsection{Нормалізована взаємокореляція}
        Нормалізована взаємокореляція шукається як взаємокореляція між сигналами, що нормалізовані.
        Як правило, під нормалізацією розуміється віднімання від сигналу середнього значення на проміжку й ділення на
        середньо"=квадратичне відхілення на цьому проміжку.

        Для дискретного сигналу нормалізована взаємокореляція має вигляд:
        \begin{equation}
            NCC(f,g)[ n ] \defeq \frac{1}{\left(k+1\right)} \sum\limits_{m=0}^{k} \left(\frac{f^*[m] -\bar{f_n}}{\sigma_{f_n}}\right) \left(\frac{g[ m + n ]-\bar{g}}{\sigma_g}\right),
        \end{equation}
        де $\sigma_{f_n}$ --- середньо"=квадратичне відхілення вікна вхідного сигналу $f$, $\sigma_g$ ---
        середньо"=квадратичне відхилення шаблону $g$, $\bar{f_n}$ --- середнє значення вхідного сигналу на проміжку
        вікна, а $\bar{g}$ --- середнє значення шаблону.

        Цей метод дозволяє знаходити шаблон, навіть якщо вхідний сигнал має значні \invcommas{спалахи} амплітуди.
        Також нормалізація шаблону дозволяє знаходити в сигналі шаблон із зміненої амплітудою.

        Завдяки тому, що значення цієї метрики лежать в межах від ${-1}$ до $1$, задача пошуку порогу $e_0$ значно
        спрощується.

    \subsection{Сума квадратів відстаней}
        Метод полягає в пошуку квадрату відстані між $k+1$"=вимірного вектору шаблону та обраного вікна сигналу.

        Для дискретного сигналу це буде наступна залежність:
        \begin{equation}
            \text{SSD}\left(f, g\right)[ n ] \defeq \sum\limits_{m=0}^k {\left(f[m] - g[m]\right)}^2
        \end{equation}

        Отримана функціональна залежність характеризує відстань кожного вікна вхідного сигналу до шаблону.
        Таким чином, чим менше ця відстань --- тим більше сигнал у обраному вікні схожий до шаблону.

        На відміну від методу взаємокореляцій, на ефективність методу пошуку суму квадратів відстаней не впливає
        значення амплітуди сигналу.

\section{Порівняльний аналіз методів}
\label{s:existing-compare}
    В таблиці~\ref{tab:existing} підсумовані основні характеристики розглянутих методів.
    Слід зазначити, що підрахування взаємокореляції можна значно прискорити завдяки (швидкому) перетворенню Фур’є
    (FFT).

    \stepcounter{tablecount}
    \begin{table}{|p{0.1\textwidth}|p{0.2\textwidth}|p{0.3\textwidth}|p{0.29\textwidth}|}
        {Порівняння існуючих методів}{tab:existing}
        {\hline
            &
Складність обчислень &
Стійкість до зміни амплітуди шаблону &
Стійкість до артефактів вхідного сигналу\\
        \hline}

        CC & $O(k n)$  & $-$ & $-$\\
        NCC & $O(2 k n)$ & $+$ & $+$\\
        SSD & $O(k n)$ & $-$ & $+$\\
    \end{table}

    Як можна побачити, найпоширеніші методи, що застосовуються для сигналів загального виду, є дуже чутливими до
    зміни шаблону: тільки метод нормалізованої взаємокореляції дозволяє коректно знаходити шаблон із лінійно"=зміненою
    амплітудою.
    \todo[inline]{Вставити посилання}
    Експерименти показали, що ці методи не знаходять шаблон в сигналі, якщо амплітуда шаблону змінювалася не сталим
    чином.

\section{Висновки}
    \todo[inline]{Обґрунтування обраного методу}

% vim: spelllang=uk,en spell filetype=tex

\chapter{Метод Кунченка}

Відомо, що поліном Тейлора є одним з найкращих способів апроксимації функції на деякому інтервалі, проте задля
використання цього методу необхідно забезпечити $n$-раз гладкість функції.

Кунченком~Ю.\,П. було запропоновано використовувати для апроксимації поліном, побудований з неортонормованого базису
функцій, що були отримані шляхом певних перетворень з шаблона.

В цьому розділі буде запропонований метод пошуку шаблонів за допомогою апроксимації поліномів у просторі Кунченко.
Для цього буде розглянутий лінійній простір Кунченко, метод розкладання довільної функції в цьому просторі й
застосування цього розкладення в проблемі пошуку шаблонів.
Також будуть розглянуті методи, за допомогою яких можна поліпшити результати алгоритму.

\section{Лінійній простір Кунченка}
    Нехай $f(x)$ --- так звана породжуюча функція, що визначена на проміжку ${[a, d]}$, а $\{ \phi_v( x ) \}$ ---
    довільна множина функціональний перетворень.
    Тоді, можна породити множину функцій
    \begin{equation}
        u_v(x) = \phi_v( f( x ))
    \end{equation}

    З цієї множини можна вибрати підмножину лінійно"=незалежних функцій.
    Лінійний простір, утворений на такій множині, називатимемо лінійним простором над незалежними породженими
    функціями, або лінійний простором Кунченка (LFKu).

    На цьому просторі можна звичайним чином визначити скалярний добуток:
    \begin{equation}
        \Psi_{v,k} = (u_v( x ), u_k( x ) ) = \int\limits_a^d u_v( x ) u_k( x )\,\mathrm{d}x,
    \end{equation}
    та відстань між двома елементами:
    \begin{equation}
        \rho_{v,k}^2 = \|u_v( x ) - u_k( x )\|^2 = \int\limits_a^d\left(u_v(x) - u_k(x)\right)^2\,\mathrm{d}x
    \end{equation}

    Виберемо деяку породжену функцію $u_b(x)$, яку будемо називати основною.
    Тоді множину функцій $u_v(x), v \ne b$ будемо називати доповнювальними функціями.

    Можна утворити поліном Кунченка з доповнювальних функцій:
    \begin{equation}
        \label{eq:polynom}
        P_t( x ) = \sum^t_{k = 1, k \ne b} a_k u_k( x  )
    \end{equation}

    Відстань між основною функцією $u_b(x)$ та утвореним поліномом за визначенням буде дорівнювати:
    \begin{equation}
        \rho^2_{b,P} = \int\limits_a^d\left( u_b(x) - P_t(x)\right)^2 \, \mathrm{d}x
    \end{equation}

    Можна показати, що коефіцієнти $a_k, k \ne 0$ знаходяться з системи лінійних рівнянь:

    \begin{equation}
        \label{eq:linear-system}
        \sum^t_{k=1,k\ne b} a_k F_{v,k} = F_{v,b},\qquad v=\overline{0,t}, v \ne b,
    \end{equation}
    де
    \begin{eqnarray}
        \label{eq:centered-correlants}
        F_{v,k} &\equiv& \Psi_{v,k} - \Psi_v \Psi_k \| u_0( x ) \|^2\\
        \Psi_v &\equiv& \frac{\Psi_{v,0}}{\left\|u_0(x)\right\|^2}\\
        a_0 &=& \Psi_b - \sum\limits_{v=1, v \ne b}^r{a_v \Psi_v}
    \end{eqnarray}

    Спростив вираз, отримаємо:
    \begin{equation}
        \rho^2_{bP} = F_{b,b} - J_t,
    \end{equation}
    де $J_t$ --- інфоркуна полінома, що має вигляд:
    \begin{equation}
        J_t \equiv \sum_{v=1,v\ne b}^t a_v F_{v,b}
    \end{equation}

    Тоді мірою апроксимації поліномом будемо вважати наступну величину:
    \begin{equation}
        \label{eq:et}
        e_t = \frac{J_t}{\int\limits_a^d (u_b(x)-\Psi_b u_0( x ))^2\, \mathrm{d}x}
    \end{equation}

\section{Застосування розкладення для пошуку шаблону в сигналі}
    Візьмемо в якості породжуючої функції шаблон $f(x)$, що є визначеним на проміжку $\forall x_i \in {[a, d]}$.
    Тоді, нехай $\{u_v(x)\}$, також визначені $\forall x_i \in {[a, d]}$ --- породжені функції від шаблону, що
    утворюють лінійній простір Кунченко.

    Якщо взяти за основну функцію $u_b(x)$ вхідний сигнал на проміжку ${[a,d]}$, за допомогою~\eqref{eq:polynom} можна
    побудувати поліном наближення сигналу $P_t(x)$.

    Оскільки поліном, визначений таким чином буде являти собою наближення вхідного сигналу за допомогою певних
    перетворень з шаблону, то можна вважати величину~\eqref{eq:et} мірою схожості вхідного сигналу до шаблону на
    обраному проміжку.

    Таким чином, скомбінував метод ковзаючого вікна (розділ~\ref{ss:sliding-window}) з розкладенням сигналу в
    лінійному просторі Кунченко, можна отримати ефектограму, тобто функцію $e_t(x)$, що буде характеризувати схожість
    вхідного сигналу зі шаблоном у певному вікні.

    Оскільки ефектограма може приймати значення від $-1$ до $1$, можливо визначити поріг $e_0$, такий, що значення
    більше (або менше при умові його від’ємності) буде характеризувати відповідне вікно як таке, що містить
    модифікований шаблон.

    Цей алгоритм потребує наступні обчислення:
    \begin{itemize}
        \item $O(t^2) O(k^2)$ обчислень для підрахунку попарних центрованих корелянт~\eqref{eq:centered-correlants}
            всіх породжених функцій (кожна з яких має $k$ значень).
            Ця дія має бути виконаною лише один раз на початку пошуку, оскільки при ковзанні вікна вони не будуть
            змінюватись.
        \item $O(t) O(k)$ обчислень для підрахунку корелянт із основною функцією (вхідним сигналом)~$F_{v,b}$ на
            кожному кроці, оскільки значення основної функції змінюється із зсувом вікна.
        \item $O(t^3)$ обчислень для пошуку коренів системи лінійних рівнянь~\eqref{eq:linear-system}.
        \item $O(t) + O(k^2)$ для обрахування ефективності наближення на кожному кроці.
    \end{itemize}

\section{Методи покращення пошуку шаблонів}
    Як видно з попередніх обчислень, запропонований метод має більшу обчислювальну складність, ніж метод із
    використанням нормалізованої взаємокореляції (розділ~\ref{s:existing-compare}).

    В цьому підрозділі будуть запропоновані методи прискорення роботи алгоритму.

    \subsection{Пірамідальний пошук}
        Оскільки основною складовою кількості обчислень можна вважати $O(k^2)$, було запропоновано використання методу
        так званого пірамідального пошуку.
        \todo[inline]{Vanderbrug, G.J. and Rosenfeld, A. (1977) Two-stage Template Matching. IEEE Transactions on Computers, C-26,}
        \todo[inline]{Rosenfeld, A. and Vanderbrug, G.J. (1977) Coarse-Fine Template Matching. IEEE Transactions on Systems, Man and Cybernetics, 7, 104-107.}

        Ідея цього методу полягає в тому, що спочатку виконується певне масштабування (стиснення) шаблону та сигналу.
        Після цього виконується пошук стисненого шаблону в стисненому сигналі.
        Якщо знаходяться екстремуми, що характеризують знахідку шаблону, то на відповідній області нестисненого
        сигналу виконується пошук нестисненого шаблону.

        Таким чином, якщо сигнал та шаблон були промасштабовані із співвідношенням $1:10$, то головна складова
        кількості обчислень зміниться до $\frac{1}{100} O(k^2)$.
        Як правило, після першого етапу залишається значно менше областей, в яких потрібно уточнювати результат.

        Також варто зазначити, що можливі два підходи щодо стиснення сигналів:
        \begin{itemize}
            \item Зміна частоти дискретизації (resampling).
                В цьому випадку можна використовувати, наприклад фільтр Ланцоша, якщо орігінальна частота
                дискретизації не ділиться на нову.
                Проте значно ефективніше обирати таку нову частоту дискретизації, що ділить орігінальну частоту
                націло.
                В такому випадку можна вибрати кожне $m$-те значення сигналу та шаблону, де $m$ --- співвідношення
                частот.
                \todo[inline]{Ken Turkowski and Steve Gabriel (1990). "Filters for Common Resampling Tasks". In Andrew
                    S. Glassner. Graphics Gems I. Academic Press. pp. 147–165. ISBN 978-0-12-286165-9.}
            \item Брати середнє значення від деякого околу точки.
        \end{itemize}
    \subsection{Метод прискореного ковзання}
        Оскільки в процесі пошуку шаблону отримуємо функціональну залежність $e_t(x)$, що характеризує ступінь
        наближення вхідного сигналу певними модифікаціями шаблону, можна розглядати задачу цього пошуку як задачу
        оптимізації.

        Найбільш ефективним застосуванням методів оптимізації буде використання алгоритм пошуку інтервалів, що могуть
        містити екстремум функції ефективності наближення.
        Такий алгоритм можна зображено на малюнку~\ref{fig:svenn-algorithm}.

        \stepcounter{figurecount}
        \begin{figure}[!h]
            \centering
            \includegraphics[height=0.8\textheight]{svenn.eps}
            \caption{Блок"=схема роботи алгоритму пошуку інтервалу з екстремумом}
            \label{fig:svenn-algorithm}
        \end{figure}

        Використовуючи цей алгоритм можна знайти інтервал, що містить певний екстремум, а потів уточнювати отримане
        значення.
        Цього можна досягти як повним обчисленням всіх значень функції $e_t(x)$ на отриманому проміжку, так і
        використанням методів уточнення інтервалів, таких як методи дихотомії та золотого перетину.

        Методи дихотомії та золотого перетину дозволяють уточнити (звузити) інтервал, що містить екстремум.
        Алгоритм методу золотого перетину має наступний вид:
        \begin{enumerate}
            \item Задати
                \begin{eqnarray*}
                    x_0 &=& a\\
                    x_1 &=& b\\
                \end{eqnarray*}
            \item Обчислити
                \begin{eqnarray*}
                    x_l &=& x_1 - \frac{x_1 - x_0}{\phi}\\
                    x_r &=& x_0 + \frac{x_1 - x_0}{\phi},
                \end{eqnarray*}
                де $\phi = \frac{1 + \sqrt{ 5 }}{2}$ --- пропорція золотого перетину.
            \item\label{enum:step3} Якщо $f(x_l) < f(x_r)$, то задати $x_0 = x_l$, інакше $x_1 = x_r$.
            \item Перевірити умови зупинки (наприклад, $\|x_1 - x_0\| < \varepsilon$).
                Якщо умова не виконується, то перейти до кроку~\ref{enum:step3}, інакше повернути інтервал ${[x_0,
                    x_1]}$ або $\frac{x_0+x_1}{2}$ в залежності від потреб задачі.
        \end{enumerate}

        \todo[inline]{Обмеження на використання}


\section{Висновки}
В цьому розділі був визначений лінійній простір над породжуючою функцією, або лінійний простір Кунченко.
Було показано, як використовуючи наближення функції елементами базису цього простору можна вирішити задачу пошуку
шаблону в цифровому сигналі.
Обраний метод має більш нечітку оцінку присутності шаблону в сигналі, ніж розглянуті в попередньому розділі.
Також були розглянуті методики, завдяки яким можна прискорити роботу методу.

% vim: spelllang=uk,en spell filetype=tex

\chapter{Проектування програмних засобів}
В цьому розділі будуть сформульовані вимоги до програмного середовища, в якому необхідно реалізувати метод пошуку
шаблонів за допомогою поліномів Кунченко.
Будуть наведені можливі варіанти вибору й обґрунтований вибір.
Буде спроектована архітектура програмної реалізації та наведений алгоритм роботи цієї реалізації.

\section{Вибір програмного середовища}
    Для полегшення розробки та забезпечення високої якості роботи програмних засобів, було вирішено використати один з
    існуючих математичних пакетів для роботи з цифровими сигналами.

    Серед інших, розглядалися наступні математичні пакети: Mathworks\textsuperscript{\textregistered}
    MATLAB\textsuperscript{\textregistered} 2014b, GNU Octave, Python+Numpy.
    Порівняльна характеристика цих пакетів наведена в таблиці~\ref{tab:math}.

    \stepcounter{tablecount}
    \begin{table}
        {|p{0.1\textwidth}|C{0.3\textwidth}|C{0.27\textwidth}|p{0.15\textwidth}|}
        {Порівняльна характеристика математичних пакетів}
        {tab:math}
        {\hline
            & Кросплатформенність & Мінімальний розмір & Ліцензія\\
            \hline}
        Matlab & + & 1200 МіБ & Propr.\\
        Octave & + & 650 МіБ  & GPL\\
        Numpy  & + & 112 МіБ  & BSD-new\\
    \end{table}

    Усі ці середовища зарекомендували себе як якісне математичне забезпечення для обробки сигналів.

    Matlab є найстарішим математичним пакетом із пакетів, що розглядалися.
    Для цієї середи було написано безліч модулів для обробки сигналів і не тільки.
    Також, із усіх пакетів, що розглядалися, тільки в цьому пакеті є вбудована можливість апаратного прискорення
    обчислень на відеочіпі, що дає значний приріст продуктивності при обробці цифрових сигналів.

    Середовище Octave було розроблено як альтернатива Matlab.
    Воно підтримує синтаксис вбудованої мови Matlab, і саме тому безліч модулів, що були написані для Matlab, працюють
    і на Octave.
    Цей пакет є програмним забезпеченням із відкритим кодом, що є значною перевагою над Matlab.
    Також варто зазначити, що цей пакет розвивається набагато швидше за Matlab, а тому, можливо, в майбутньому в ньому
    будуть функції, що зараз є лише в Matlab (наприклад, прискорення обчислень за допомогою апаратного прискорення).

    Numpy, на відміну від Matlab та Octave, є набором модулів/бібліотек для мови програмування Python.
    З усіх систем, Numpy розвивається найактивніше.
    До переваг Numpy можна віднести те, що програми на Numpy --- це програми на мові Python, яка є набагато
    поширенішою за мови Matlab/Octave.
    Також ліцензія, за якою розповсюджується Numpy є більш вільною за GPL (за якою розповсюджується Octave), що
    дозволяє використовувати це середовище для комерційної розробки.

    Для початкового проектування й дослідження було обрано середовище Matlab, оскільки знайти необхідний модуль для
    цього середовища значно легше, а можливість використовувати апаратне прискорення також є невід’ємною
    характеристикою.
\section{Проектування архітектури програмної реалізації}
    Оскільки було обрано математичне середовище Matlab, програмна реалізація методу пошуку шаблонів за допомогою
    розкладення функції поліномами Кунченка складається із наступних функцій"=скриптів:
    \begin{itemize}
        \item \verb'match' --- функція"=інтерфейс, точка входу до модуля пошуку шаблонів.
            Отримує параметрами цифровий сигнал, шаблон, функціональні перетворення, за допомогою яких будується
            лінійний простір над породжуючою функцією.
            Саме в цій функції виконується \invcommas{ковзання} вікна.
            Ця функція, в свою чергу, викликає \verb'buildGeneratedFunctionsSystem', \verb'createIterator',
            \verb'approximate'.
        \item \verb'buildGeneratedFunctionsSystem' --- функція, що будує базиси лінійного простору Кунченко та
            обчислює центровані скалярні добутки між ними.
        \item \verb'createIterator' --- створює вікно й дозволяє рухати це вікно заданим чином.
        \item \verb'approximate' --- обчислює апроксимацію обраного вікна вхідного сигналу базисом лінійного простору
            Кунченко.
            В цій функції будується система лінійних рівнянь, з якої знаходяться коефіцієнти при базисних функціях
            лінійного простору, рахується коефіцієнт наближення вхідного сигналу за допомогою породжених функцій.
    \end{itemize}

    Таким чином було створено наступні рівні абстракції:
    \begin{itemize}
        \item Базис лінійного простору, що містить наступні поля:
            \begin{itemize}
                \item породжуюча функція (шаблон);
                \item породжені функції (функціональні перетворення шаблону);
                \item посилання на функцію скалярного добутку;
                \item наперед підраховані значення скалярного добутку для усіх пар породжених функцій.
            \end{itemize}

            Посилання на функцію скалярного добутку в цій структурі дозволяє змінювати метод пошуку корелянти між
            двома функціями.
            За замовчанням це є чисельне інтегрування методом трапецій від попарного добутку значень аргументів.
        \item Ітератор в сигналі, який дозволяє певним чином обходити довільний сигнал.
            Змінюючи реалізацію цього ітератору можна змінювати вигляд вікна вхідного сигналу, змінювати шаг між
            вікнами, тощо.
    \end{itemize}

\section{Алгоритм роботи програми}
    \stepcounter{figurecount}
    \begin{figure}[!h]
        \centering
        \includegraphics[height=0.82\textheight]{Drawing2.eps}
        \caption{Блок"=схема роботи алгоритму пошуку шаблону}
        \label{fig:block_scheme_monitor}
    \end{figure}

\section{Висновки}
    В цьому розділі було обґрунтовано вибір середовища Matlab як середовища з великою кількістю модулів для обробки
    цифрових сигналів.

    Завдяки абстрагуванню спроектована система має достатню гнучкість, щоб змінювати метод пошуку наступного вікна та
    обчислення скалярного добутку.
    Завдяки цьому отримана система дозволяє шукати шаблон не лише в одновимірному сигналі, а й в двовимірному.

    Також можна зазначити, що після дослідження методу пошуку шаблонів за допомогою розкладення функції в лінійному
    просторі Кунченка, систему можна буде перевести в середовище Numpy, щоб можна було користуватися багатьма
    системними бібліотеками, що розроблені для мови програмування Python.
    Це дозволить програмному продукту мати, наприклад, графічний інтерфейс користувача.

    Також має сенс мати сумісність з математичним середовищем Octave, адже останнє, на відміну від Matlab, є
    безкоштовним ПЗ з відкритим сирцевим кодом.

% vim: spelllang=uk,en spell filetype=tex


\chapter{Математичне моделювання}


% vim: spelllang=uk,en spell


\conclusions{}
В магістерській дисертації розглядалася задача пошуку шаблонів у цифрових сигналах.
В ході дослідження було проаналізовані існуючі методу для пошуку шаблонів і порівняні їх обмеження.

Під час огляду предметної області були розглянуті такі поняття, як цифровий сигнал, кросс"=кореляція, сума квадратів
відстаней, метод ковзного вікна та віконні функції.
Розглянуті віконні функції було використано при тестуванні в методі ковзного вікна.

Був запропонований метод пошуку шаблонів, що базується на розкладанні функції в базисі лінійного функціонального
простору Кунченка.
Відмінною особливістю цього методу є можливість пошуку шаблонів, що були значно змінені в сигналі.

Також були запропоновані шляхи покращення роботи методу на основі поліномів Кунченка.
Пірамідальний пошук дозволяє значно прискорити роботу методу, оскільки він дозволяє відфільтрувати значну частину
сигналу як таку, що не містить шаблону, за досить короткий проміжок часу.
Використання методів оптимізації також дозволяють значно прискорити роботу алгоритму.

Був розроблений та оптимізований програмний комплекс в математичному середовищі MathWorks MATLAB 2014b.

Під час статистичного експерименту було порівняно ефективність роботи обраного методу з існуючими методами як на
штучно"=згенерованому сигналі, так і для аудіозапису мовлення.
Як результат цього тестування можна зазначити, що обраний метод дозволяє отримати аналогічну ступінь розпізнавання на
штучно"=згенерованих сигналах і значно вищу якість на аудіозаписах мовлення.

Для пошуку шаблону в аудіозаписі з мовленням було використано визначення коротко"=часової енергії й віконних функцій.
Було проаналізовано поведінку алгоритму в залежності від обраних параметрів методу.

Для подальшого розвитку слід зазначити можливість дослідження методів інтелектуального вибору параметрів алгоритму в
залежності від характеристик сигналу, апаратне прискорення для більш швидкої роботи алгоритму та апробація методу на
даних з іншої предметної області.

\end{document}
% vim: spelllang=uk,en spell
